% Part: first-order-logic
% Chapter: introduction
% Section: models-theories

\documentclass[../../../include/open-logic-section]{subfiles}

\begin{document}

\olfileid{fol}{int}{mod}

\olsection{模型与理论}

一旦定义了一阶逻辑的语法和语义，我们就可以着手研究结构的性质和各种语义概念了。
我们也可以定义推导系统并研究它们。
对于一个语句集，我们可以提问：什么样的结构能使这个集合中的所有语句为真？
给定语句集~$\Gamma$，满足其中所有语句的结构~$\Struct{M}$ 就称为~$\Gamma$ 的一个\emph{模型}。
我们可能会从~$\Gamma$ 出发，试图找到它的模型——它们是什么样的？论域必须多大或多小？
但我们也可能会从单个结构或一类结构出发，然后提问：哪些语句在它们中为真？
是否存在一些语句能够\emph{刻画}这些结构，即这些语句在且仅在这些结构中为真？
这类问题正是\emph{模型论}的研究领域。
它们也是\emph{公理化方法}的基础：用一组语句——即一个理论的公理——来描述一类结构。
这种方法之所以可行，是因为我们观察到：在一阶逻辑中，恰恰是公理所语义后承的语句，才在公理的所有模型中为真。

举一个非常简单的例子，考虑预序。预序是某个集合~$A$ 上的一种关系~$R$，它既是自反的又是传递的。
一个带有一个二元关系 $R \subseteq A \times A$ 的集合~$A$，恰好足以给出仅含一个二元谓词符号~$\Obj P$ 的一阶语言的一个结构：我们令 $\Domain{M} = A$ 且 $\Assign{\Obj P}{M} = R$。
既然 $R$~是预序，它是自反且传递的，我们可以找到一阶逻辑的一个语句集~$\Gamma$ 来表达这一点：
\begin{align*}
  & \lforall[\Obj v_0][\Atom{\Obj P}{\Obj v_0,\Obj v_0}]\\
  & \lforall[\Obj v_0][\lforall[\Obj v_1][\lforall[\Obj v_2][((\Atom{\Obj P}{\Obj v_0,\Obj v_1} \land \Atom{\Obj P}{\Obj v_1,\Obj v_2}) \lif \Atom{\Obj P}{\Obj v_0,\Obj v_2})]]]
\end{align*}
这些语句正是“对任意~$x$，$Rxx$”（$R$ 是自反的）和“每当 $Rxy$ 且 $Ryz$，则也有 $Rxz$”（$R$ 是传递的）的符号化。
我们看到，结构~$\Struct{M}$ 是这两个语句组成的集合~$\Gamma$ 的模型，当且仅当 $R$（即 $\Assign{\Obj P}{M}$）是 $A$（即 $\Domain{M}$）上的一个预序。
换言之，$\Gamma$ 的模型恰好就是所有的预序。
所有预序的任何性质，只要能用以~$\Obj P$ 为唯一谓词符号的一阶语言表达（如上文的自反性和传递性），都是~$\Gamma$ 中这两个语句的语义后承，反之亦然。
因此，凡是我们能对~$\Gamma$ 的模型证明的结论，也就对所有预序得到了证明。

对于任何特定的理论和模型类（例如 $\Gamma$ 和所有预序），都会存在有趣的问题：在相应的一阶语言中，什么可以表达，什么不可以表达。
有一些结构的性质对于所有语言和模型类都是值得关注的，即那些关于论域大小的性质。
例如，总可以表达论域恰好包含 $n$~个元素，其中~$n \in \PosInt$。
也可以用无穷多条语句的集合来表达论域是无穷的。
但是，无法表达论域是有穷的，或者论域是不可数的。
这些关于一阶语言局限性的结果，是紧致性定理和 L\"owenheim（勒文海姆）--Skolem（斯科伦）定理的推论。

\end{document}